Large language models, such as GPT-4~\citep{openai2023gpt4}, GPT-3~\citep{brown2020language}, PaLM~\citep{chowdhery2022palm}, PaLM-2~\citep{anil2023palm}, BLOOM~\citep{workshop2023bloom}, OPT~\citep{zhang2022opt}, LLaMA~\citep{touvron2023llama}, Llama-2~\citep{touvron2023llama2}, and many others~\citep{radford2019language, gpt-j, gpt-neo, smith2022using, hoffmann2022training, penedo2023refinedweb} have shown remarkable performance on natural language processing (NLP) tasks. Following the pretraining phase, additional finetuning enables models (e.g., FLAN~\citep{wei2022finetuned}, FLAN-T5~\citep{chung2022scaling}, InstructGPT~\citep{ouyang2022training}) to better align with human instructions to complete tasks. Moreover, instruction-tuning has enabled LLMs (e.g., GPT-3.5 Turbo~\citep{openaichatgpt}, Llama-2-chat~\citep{touvron2023llama2}, Self-Instruct~\citep{wang2023selfinstruct}, Alpaca~\citep{alpaca2023}, Vicuna~\citep{vicuna2023}, Koala~\citep{koala2023}) to better engage with users through dialogue. \zeroshotstart\ \method\ builds on instruction-following language models, enabling better zero-shot reasoning abilities through the use of agent instructions.

Language agents~\citep{yao2023react, shinn2023reflexion, xu2023expertprompting, park2023generative, zhou2023agents, andreas-2022-language, wang2023survey, xi2023rise, sumers2023cognitive, chan2023chateval} recently emerged due to the task planning capabilities of LLMs. Given a task often demonstrated in natural language, these agents aim to complete the task directly. 
We base our agent on ReAct~\citep{yao2023react}, which uses structured ``Thought, Action, Observation" steps to guide a model to accomplish a goal.
Unlike existing agents that aim to solve a task directly, our agent generates task-specific instructions on how to complete the given task, decoupling the agent planning and reasoning steps. Besides reaching competitive performance with using agents directly, our design turns out to be more cost-effective.

Finetuning is an effective method for generating higher-quality responses from LLMs on downstream tasks~\citep{liu2019roberta, howard2018universal}. As model scale increases, finetuning becomes less practical, which lightweight tuning methods (such as prefix tuning~\citep{li2021prefixtuning}, prompt learning~\citep{lester2021power, liu2023pre}, LoRA~\citep{hu2021lora}) have tried to solve. Even with such methods, prompting techniques have gained attention as an even cheaper alternative, such as chain of thought and zero-shot chain of thought prompting~\citep{wei2023chainofthought, kojima2023large}, the latter of which we focus on in our work.
Few-shot learning, which involves providing a few examples demonstrating the task before prompting the models during inference, is often effective on a range of tasks~\citep{brown2020language, dong2023survey}. Chain of thought (CoT) prompting~\citep{wei2023chainofthought} involves generating a series of intermediate reasoning steps, which can dramatically increase the performance of LLMs on complex reasoning tasks. While this reasoning behavior is traditionally learned from few-shot demonstrations, \citet{kojima2023large} extends CoT prompting to the zero-shot setting, where no examples are used. More recently, new approaches such as self-consistency~\citep{wang2023selfconsistency}, plan-and-solve prompting~\citep{wang2023planandsolve}, tree of thought~\citep{yao2023tree}, and graph of thought~\citep{besta2023graph} have further improved the reasoning quality. \zeroshotstart\ \method\ focuses on the zero-shot setup. It generalizes the reasoning abilities of LLMs to more tasks by utilizing task-specific instructions generated from our agent to better align the reasoning process with a particular task.

NLP benchmarking datasets provide a standardized interface to evaluate LLMs on specific downstream tasks. Common benchmarks (e.g., HELM~\citep{liang2022holistic}, MMLU~\citep{hendrycks2021measuring}, and many others~\citep{shen2024measuring, yuan2024measuring, zhong2023agieval, srivastava2023imitation, suzgun2022challenging, chen2021evaluating, wang2019glue, rajpurkar-etal-2016-squad}) have become part of the standard evaluation of LLMs. We benchmark our method on 29 datasets, consisting of the core scenarios from HELM~\citep{liang2022holistic} and the reasoning datasets from~\citet{kojima2023large}, and include results for comparison purposes.
For reasoning tasks, our method significantly outperforms existing zero-shot approaches~\citep{brown2020language,kojima2023large}.
Besides reasoning tasks, our method generalizes to general language understanding tasks including generation and classification.